%----------------------------------------------------------------------------------------
% BREUKEN: GETALLEN EN LETTERVORMEN. Alle opgaven zijn eigen werk.
%
% Een hoofdstuk kernvaardigheden (ontwerpregels, deel Kernvaardigheden): beknopte theorie, daarna drill.
% Labels van dit hoofdstuk beginnen met h4-.
%----------------------------------------------------------------------------------------

\mychapter{4}{Breuken}

\lkr{Herhaling van rekenen met breuken (1ste graad) en uitbreiding naar breuken met letters (rationale vormen). Nodig voor vergelijkingen met breuken (hoofdstuk 5), de functie $c/x$ (Functies), en later voor rationale functies in de 3de graad. Veronderstelt hoofdstuk 3 (ontbinden). Wijs bij bundel 3 op de bestaansvoorwaarden: de vereenvoudigde breuk geldt alleen waar de oorspronkelijke bestaat.}

\section{Wat je moet kennen}

\begin{theorie}
\begin{minipage}[t]{.5\linewidth}
\vspace{0pt}
$\formule{\dfrac ab + \dfrac cd = \dfrac{ad + bc}{bd}}$ \qquad $\formule{\dfrac ab \cdot \dfrac cd = \dfrac{ac}{bd}}$

\vspace{2mm}
$\formule{\dfrac ab : \dfrac cd = \dfrac ab \cdot \dfrac dc}$ \qquad $\formule{\dfrac{a \cdot c}{b \cdot c} = \dfrac ab}$ \ ($c \neq 0$)
\end{minipage}\hfill
\begin{minipage}[t]{.46\linewidth}
\vspace{0pt}
Een breuk bestaat niet als de noemer $0$ is: dat geeft de \textbf{\textcolor{accent}{bestaansvoorwaarden}}.

\vspace{1mm}
Je vereenvoudigt door teller en noemer te \emph{ontbinden} en dezelfde \emph{factor} te schrappen. Optellen doe je op de kleinste gemeenschappelijke noemer.
\end{minipage}
\end{theorie}

\textbf{Klassieke fouten.} Schrappen in een som: $\dfrac{x + 3}{3} \neq x$, want $3$ is geen factor van de teller. Noemers optellen: $\dfrac1a + \dfrac1b \neq \dfrac{2}{a + b}$. Het minteken voor een breuk vergeten bij de hele teller: $-\dfrac{x - 1}{2} = \dfrac{-x + 1}{2}$.

\voorbeeld{}{Vereenvoudig $\dfrac{x^2 - 9}{x^2 + 3x}$ en geef de bestaansvoorwaarden.\label{vb:h4-vereenvoudig}}{%
\begin{stappen}
\frac{(x - 3)(x + 3)}{x(x + 3)} = \frac{x - 3}{x} & Ontbinden, en de factor $x + 3$ schrappen. \\
x \neq 0 \text{ en } x \neq -3 & Uit de oorspronkelijke noemer $x(x + 3)$.
\end{stappen}}

\voorbeeld{}{Werk uit: $\dfrac{2}{x - 1} - \dfrac{3}{x + 1}$.\label{vb:h4-optellen}}{%
\begin{stappen}
\frac{2(x + 1) - 3(x - 1)}{(x - 1)(x + 1)} = \frac{5 - x}{(x - 1)(x + 1)} & Gelijke noemer $(x - 1)(x + 1)$; haakjes rond de tellers.
\end{stappen}}

%========================================================================================
\opdrachtenkop

\begin{oefeningen}
\needspace{10\baselineskip}
\oef{h4-o1}

\needspace{10\baselineskip}
\oef{h4-o2}

\needspace{10\baselineskip}
\oef{h4-o3}

\needspace{10\baselineskip}
\oef{h4-o4}

\needspace{10\baselineskip}
\oef{h4-o5}

\needspace{10\baselineskip}
\oef{h4-o6}

\needspace{10\baselineskip}
\oef{h4-o7}

\needspace{10\baselineskip}
\oef{h4-o8}

\needspace{10\baselineskip}
\oef{h4-o9}

\needspace{10\baselineskip}
\oef{h4-o10}
\end{oefeningen}

%----------------------------------------------------------------------------------------
\needspace{14\baselineskip}
\opdrachtsectie{Test jezelf}

\begin{oefeningen}
\oef{h4-t1}

\oef{h4-t2}

\oef{h4-t3}

\oef{h4-t4}
\end{oefeningen}

\herhaalschema{%
  \fzv{oef:h4-t1}{oef:h4-o1}%
  \fv{oef:h4-t2}{vb:h4-vereenvoudig}{oef:h4-o3}%
  \fv{oef:h4-t3}{vb:h4-vereenvoudig}{oef:h4-o4}%
  \fv{oef:h4-t4}{vb:h4-optellen}{oef:h4-o5}}

\ifzelfstudie
\section*{Maak je eigen samenvatting}
Sluit dit boek en schrijf in je schrift: hoe je breuken optelt, vermenigvuldigt en deelt, wat bestaansvoorwaarden zijn, hoe je een breuk met letters vereenvoudigt, en wanneer je \emph{niet} mag schrappen. Vergelijk daarna met het kader.
\fi
